% !TEX root = ../main.tex
\corrigesChapitre{1}{Polynômes du second degré}
\corrigesSouples

\begin{blocLivre}{cor-1-1}
\begin{corrige}{1}{Reconnaître un trinôme}
\methode{développer pour identifier les coefficients}
Les fonctions du second degré sont $3x^2-2x+5$ et $(x-2)^2=x^2-4x+4$.
Leurs coefficients sont respectivement $(3;-2;5)$ et $(1;-4;4)$.
La fonction $4x-1$ est de degré 1 ; $x^3+x$ est de degré 3.
\end{corrige}
\end{blocLivre}
\begin{blocLivre}{cor-1-2}
\begin{corrige}{2}{Développer ou factoriser}
\demarche{Choisir entre double distributivité, identité remarquable et facteur commun.}
\begin{questions}
\item Par double distributivité, $(x-3)(x+2)=x^2+2x-3x-6=\res{x^2-x-6}$.
Puis $(x-1)^2=x^2-2x+1$, donc $2(x-1)^2-8=2x^2-4x+2-8=\res{2x^2-4x-6}$.
\item Une différence de carrés donne $x^2-25=x^2-5^2=\res{(x-5)(x+5)}$.
Le facteur commun $3x$ donne $3x^2-12x=\res{3x(x-4)}$. Développer ces produits permet de les contrôler.
\end{questions}

\end{corrige}
\end{blocLivre}
\begin{blocLivre}{cor-1-3}
\begin{corrige}{3}{Lire la forme canonique}
\demarche{Lire le sommet dans le carré, puis isoler ce carré pour les racines.}
\begin{questions}
\item Le carré est minimal en $x=3$ : sommet $\res{S(3;-8)}$, axe $x=3$, minimum $-8$.
\item $f(0)=18-8=10$. $f(x)=0\iff(x-3)^2=4\iff x-3=\pm2$,
donc $\res{x=1\text{ ou }x=5}$.
\end{questions}
\end{corrige}
\end{blocLivre}
\begin{blocLivre}{cor-1-4}
\begin{corrige}{4}{Calculer un discriminant}
\demarche{Identifier les trois coefficients avant de calculer le discriminant.}
\begin{questions}
\item Ici $(a;b;c)=(1;-5;6)$ et $\Delta=(-5)^2-4\times1\times6=1>0$.
Les solutions sont $(5-1)/2=2$ et $(5+1)/2=3$, donc $\res{S=\{2;3\}}$.
\item Ici $\Delta=4^2-4\times1\times4=0$. La racine double vaut $-4/(2\times1)=-2$ : $\res{S=\{-2\}}$.
On retrouve $(x+2)^2=0$.
\item Ici $(a;b;c)=(2;1;2)$ et $\Delta=1^2-4\times2\times2=-15<0$.
Le trinôme ne s’annule pas sur $\R$ : $\res{S=\varnothing}$.
\end{questions}

\end{corrige}
\end{blocLivre}
\begin{blocLivre}{cor-1-5}
\begin{corrige}{5}{Signe d'une forme factorisée}
\demarche{Étudier les deux facteurs et tenir compte du coefficient négatif.}
Les facteurs s’annulent en $-1$ et 4. Chacun est négatif à gauche de sa racine et positif à droite.
On multiplie leurs signes, puis on inverse le signe obtenu à cause du coefficient $-2$.
Le tableau ci-dessous donne un signe positif entre les racines et négatif à l’extérieur.
Ainsi $\res{P(x)\geq0\iff x\in[-1;4]}$ et
$\res{P(x)<0\iff x\in\mathopen{]}-\infty;-1[\cup]4;+\infty[}$.
Les racines sont incluses pour l’inégalité large, et exclues pour l’inégalité stricte.

\label{figlivre-r01-05-signes-cor}
\figureLivre[corrige]{r01-05-signes}

\end{corrige}
\end{blocLivre}
\begin{blocLivre}{cor-1-6}
\begin{corrige}{6}{Choisir la bonne forme}
\demarche{Choisir la forme qui rend chaque équation ou inéquation la plus courte.}
\begin{questions}
\item $(x-3)^2-4=x^2-6x+5=(x-1)(x-5)$ en développant.
\item $f(x)=5\iff x(x-6)=0$ : $\res{x=0\text{ ou }6}$.
Le trinôme est négatif ou nul entre ses racines : $\res{x\in[1;5]}$ pour $f(x)\leq0$.
\item La forme canonique donne $(x-3)^2\leq1$, donc $-1\leq x-3\leq1$ :
$\res{x\in[2;4]}$.
\end{questions}
\end{corrige}
\end{blocLivre}
\begin{blocLivre}{cor-1-7}
\begin{corrige}{7}{Deux paraboles}
\demarche{Comparer les courbes en étudiant la différence des deux fonctions.}
\begin{questions}
\item Le dessin suggère des intersections pour $x=0$ et $x=3$.
\item $f(x)-g(x)=2x^2-6x=2x(x-3)$.
Les intersections sont donc $\res{(0;3)\text{ et }(3;0)}$.
La différence est négative ou nulle entre 0 et 3, donc $\res{f\leq g\iff x\in[0;3]}$.
\end{questions}
\methode{le dessin suggère ; la factorisation démontre}
\end{corrige}
\end{blocLivre}
\begin{blocLivre}{cor-1-8}
\begin{corrige}{8}{Un quotient et une valeur interdite}
\demarche{Exclure le zéro du dénominateur avant toute étude de signe.}
Le domaine est $\R\setminus\{1\}$ et $x^2-5x+6=(x-2)(x-3)$.
\begin{center}\begin{tableau}
\otTab[2.1]{$x$/0.8,$x-2$/0.7,$x-3$/0.7,$x-1$/0.7,$Q(x)$/0.8}{$-\infty$,$1$,$2$,$3$,$+\infty$}
\tkzTabLine{,-,t,-,z,+,t,+,}
\tkzTabLine{,-,t,-,t,-,z,+,}
\tkzTabLine{,-,z,+,t,+,t,+,}
\tkzTabLine{,-,d,+,z,-,z,+,}
\end{tableau}\end{center}
Ainsi $\res{S=]1;2]\cup[3;+\infty[}$. Le point 1 est exclu ; 2 et 3 sont inclus.
\end{corrige}
\end{blocLivre}
\begin{blocLivre}{cor-1-9}
\begin{corrige}{9}{Retrouver une fonction}
\demarche{Utiliser les racines pour factoriser, puis le point donné pour fixer le coefficient.}
\begin{questions}
\item Les deux racines donnent $f(x)=a(x+2)(x-4)$. Avec $f(0)=-8$,
$-8a=-8$, donc $\res{a=1}$.
\item $f(x)=x^2-2x-8=(x-1)^2-9$. Le sommet est $\res{S(1;-9)}$.
\item $f(x)=-5\iff(x-1)^2=4$, donc $\res{x=-1\text{ ou }3}$.
\end{questions}
\label{figlivre-r01-09-parabole-cor}
\figureLivre[corrige]{r01-09-parabole}

\end{corrige}
\end{blocLivre}
\begin{blocLivre}{cor-1-10}
\begin{corrige}{10}{Somme et produit des racines}
\demarche{Développer le produit associé aux deux racines pour relier somme et produit.}
Si $u$ et $v$ sont les deux nombres cherchés, $(x-u)(x-v)=x^2-(u+v)x+uv$.
\begin{questions}
\item La somme 7 et le produit 10 conduisent à $x^2-7x+10=0$.
Comme $x^2-7x+10=(x-2)(x-5)$, les nombres sont $\res{2\text{ et }5}$.
\item La somme 4 et le produit 5 conduisent à $x^2-4x+5=0$.
Son discriminant vaut $16-20=-4<0$ : il n’existe aucun couple de réels qui convienne.
\item Pour $x^2+x-12$, la somme des racines doit être $-1$ et leur produit $-12$.
Les nombres $-4$ et 3 conviennent, donc $\res{x^2+x-12=(x+4)(x-3)}$.
\end{questions}

\end{corrige}
\end{blocLivre}
\begin{blocLivre}{cor-1-11}
\begin{corrige}{11}{Discuter selon un paramètre}
\demarche{Comparer le minimum du trinôme à zéro, en gardant le paramètre fixé.}
\begin{questions}
\item $x^2-4x+m=\res{(x-2)^2+m-4}$.
\item Résoudre revient à $(x-2)^2=4-m$. Si $m<4$, deux solutions
$2\pm\sqrt{4-m}$ ; si $m=4$, une solution double 2 ; si $m>4$, aucune solution réelle.
\item Pour $m=3$ : $\{1;3\}$ ; pour $m=4$ : $\{2\}$ ; pour $m=5$ : $\varnothing$.
\item Le minimum, atteint en 2, vaut $m-4$. Il est positif ou nul si et seulement si $\res{m\geq4}$.
\end{questions}
\end{corrige}
\end{blocLivre}
\begin{blocLivre}{cor-1-12}
\begin{corrige}{12}{Un jardin contre un mur}
\demarche{Traduire la longueur de clôture avant de chercher le maximum de l’aire.}
\begin{questions}
\item La longueur restante est $40-2x>0$ ; avec $x>0$, on obtient $0<x<20$.
L'aire vaut $A(x)=x(40-2x)=\res{-2x^2+40x}$.
\item $A(x)=-2(x-10)^2+200$. Le carré étant positif ou nul, le maximum est
$\res{200\text{ m}^2}$ pour $x=10$ m ; la longueur vaut 20 m.
\item $A(x)\geq150\iff(x-10)^2\leq25$,
donc $\res{x\in[5;15]}$ mètres, valeurs toutes admissibles.
\end{questions}
\end{corrige}
\end{blocLivre}
\begin{blocLivre}{cor-1-13}
\begin{corrige}{13}{Une vente rentable}
\demarche{Calculer recette moins coût, puis limiter la conclusion au domaine commercial.}
\begin{questions}
\item La recette est le prix multiplié par la quantité :
$R(q)=q(50-0{,}2q)=-0{,}2q^2+50q$.
Le bénéfice est $B=R-C$, donc $\res{B(q)=-0{,}2q^2+40q-500}$.
\item En développant $-0{,}2(q-100)^2+1500$, on retrouve $B$.
Le maximum sur $[0;100]$ vaut $\res{1500\text{ euros}}$, atteint pour $q=100$.
\item $B(q)\geq0\iff(q-100)^2\leq7500$.
Les racines sont $100\pm50\sqrt3$, soit environ $13{,}40$ et $186{,}60$.
En tenant compte du domaine, $q\in[100-50\sqrt3;100]$.
Les quantités entières rentables sont donc $\res{14\leq q\leq100}$.
\end{questions}
\begin{attention}
Le seuil supérieur calculé dépasse le domaine du modèle. On ne prolonge pas
la conclusion économique au-delà de 100 articles.
\end{attention}
\end{corrige}
\end{blocLivre}
\begin{blocLivre}{cor-1-14}
\begin{corrige}{14}{Un ballon au-dessus d'un filet}
\demarche{Étudier la hauteur sur les temps physiquement admissibles.}
\begin{questions}
\item $h(0)=1$ m et $h(t)=-5(t-1)^2+6$ : maximum $\res{6\text{ m}}$ à $t=1$ s.
\item Les solutions sont $1\pm\sqrt{30}/5$. Seule la positive est admissible :
$\res{t_{\text{sol}}=1+\sqrt{30}/5\approx2{,}10\text{ s}}$.
\item $h(t)\geq4\iff(t-1)^2\leq2/5$, donc
$\res{t\in[1-\sqrt{10}/5;1+\sqrt{10}/5]}$, environ $[0{,}37;1{,}63]$ s.
\item $h(1{,}5)=-11{,}25+15+1=4{,}75$ m.
Le ballon franchit le filet avec une marge verticale de $\res{0{,}75\text{ m}}$.
\end{questions}
\label{figlivre-r01-14-ballon-cor}
\figureLivre[corrige]{r01-14-ballon}

\end{corrige}
\end{blocLivre}
\begin{blocLivre}{cor-1-15}
\begin{corrige}{15}{Une famille de paraboles}
\demarche{Étudier le signe du discriminant comme une expression du paramètre.}
\begin{questions}
\item $\Delta=(-2m)^2-4(m+2)=\res{4(m-2)(m+1)}$.
\item Il y a deux racines distinctes si $m<-1$ ou $m>2$ ; une double si
$m=-1$ ou $m=2$ ; aucune si $-1<m<2$.
\item $P_m(1/2)=1/4-m+m+2=9/4$ : le point commun est $\res{A(1/2;9/4)}$.
\item $P_3(x)=x^2-6x+5=(x-1)(x-5)$. Ainsi $P_3(x)\leq0$ pour
$\res{x\in[1;5]}$ et son sommet est $\res{S(3;-4)}$.
\item Le coefficient dominant vaut 1. Le trinôme est strictement positif
partout exactement lorsqu'il n'a aucune racine réelle : $\res{-1<m<2}$.
\end{questions}
\label{figlivre-r01-15-parametre-cor}
\figureLivre[corrige]{r01-15-parametre}

\end{corrige}
\end{blocLivre}
